\par\noindent\SourceLeaf{53}%
et generalis $h^{\text{tus}}$

\[
= \frac{1}{2^{2q+2h-2}} {}^{2q+2h-1}\mathfrak{B}^{(2q+h-1)}  .  {}^{\nu}\mathfrak{B}^{(2q+2h-1)} n^{2q+2h-1}.
\]

C) Eadem ratione lex coefficientis in cosinum multipli alicuius paris cadentis, v. g. in $\cos 2q\varphi$, eruitur. In solis enim potestatibus ipsius $\cos\varphi$ paribus, inde ab illa, cuius exponens est $2q$, eiusmodi multipli cosinus occurrit. Quare seriei, qua coefficiens ipsius $\cos 2q\varphi$ exprimitur, terminus initialis ex eo seriei fundamentalis membro, quo $\cos\varphi^{2q}$ continetur, i. e. ex ${}^{\nu}\mathfrak{B}^{(2q)} n^{2q}\cos\varphi^{2q}$, et quidem ex huius termino $q^{\text{to}}$ sumetur, fietque (cum sit $\cos\varphi^{2q}$ term. $q^{\text{tus}}$

\[
= \frac{1}{2^{2q-1}} {}^{2q}\mathfrak{B}^{(2q)} \cos 2q\varphi = \frac{1}{2^{2q-1}} \cos 2q\varphi) \qquad = \frac{1}{2^{2q-1}} {}^{\nu}\mathfrak{B}^{(2q)} n^{2q}.
\]

\noindent Generatim ex $h^{\text{to}}$ post hoc primum seriei fundamentalis membro pari, adeoque ex quantitate ${}^{\nu}\mathfrak{B}^{(2q+2h)} n^{2q+2h} \cos\varphi^{2q+2h}$ proveniet coefficientis, ad $\cos 2q\varphi$ pertinentis, pars $h^{\text{ta}}$. Est autem ipsius $\cos\varphi^{2q+2h}$, si evolvatur, terminus $q^{\text{tus}}$

\[
= \frac{1}{2^{2q+2h-1}} {}^{2q+2h}\mathfrak{B}^{(2q+h)} \cos 2q\varphi.
\]

\noindent Sic igitur provenit coefficientis, ad $\cos 2q\varphi$ pertinentis terminus post primum $h^{\text{tus}}$

\[
= \frac{1}{2^{2q+2h-1}} {}^{2q+2h}\mathfrak{B}^{(2q+h)}  .  {}^{\nu}\mathfrak{B}^{(2q+2h)} n^{2q+2h}.
\]

D) Perfectum revera est negotium propositum. Non solum enim seriei ex $(1+n\cos\varphi)^{\nu}$, dum pro quavis potestate series cosinuum multiplorum substituitur, rite ordinatae terminus initialis (cf. A), verum etiam pro utroque genere terminorum, qui in illa occurrere possunt, quorum alteri continent cosinum multipli alicuius imparis, $\cos(2q-1)\varphi$ (cf. B), alteri cosinum multipli alicuius paris $\cos 2q\varphi$ (conf. C), debiti coefficientis terminus initialis et generalis rite exhibitus est. Sed haec delineatio legis inventae contrahitur observatione prorsus eandem esse legem, quae in coefficientibus cosinuum parium, et quae in imparium valet. Quod sic perspicitur. Erat coefficientis ad $\cos(2q-1)\varphi$ pertinentis terminus initialis

\[
\frac{1}{2^{2q-2}} {}^{\nu}\mathfrak{B}^{(2q-1)} n^{2q-1},
\]

\noindent terminus generalis $h^{\text{tus}}$

\[
\frac{1}{2^{2q+2h-2}} {}^{2q+2h-1}\mathfrak{B}^{(2q+h-1)}  .  {}^{\nu}\mathfrak{B}^{(2q+2h-1)} n^{2q+2h-1}.
\]
